% =============================================================================
%  MATERIAL DE ESTUDIO: FLUJO DE PROGRAMACIÓN PARA PRINCIPIANTES
%  Diagramas Mermaid — 10 ejercicios resueltos y explicados
%  Compilable en Overleaf (XeLaTeX o pdfLaTeX)
% =============================================================================
\documentclass[10pt,a4paper]{article}

% ---------------------- Paquetes básicos ------------------------
\usepackage[utf8]{inputenc}
\usepackage[T1]{fontenc}
\usepackage[spanish]{babel}
\usepackage{geometry}
\geometry{margin=2.5cm}
\usepackage{hyperref}
\usepackage{enumitem}
\usepackage{bookmark}
\usepackage{xcolor}
\usepackage{fancyvrb}
\usepackage{listings}
\usepackage{tcolorbox}
\usepackage{mdframed}
\usepackage{parskip}

% ---------- Estilo para el diagrama Mermaid (textual) ----------
\newenvironment{mermaid}[1][]
  {%
   \begin{mdframed}[%
     linecolor=blue!50!black,
     backgroundcolor=blue!3,
     innertopmargin=6pt,
     innerbottommargin=6pt,
     leftmargin=0pt,
     rightmargin=0pt,
     skipabove=10pt,
     skipbelow=10pt,
     frametitle={Diagrama Mermaid},
     frametitlefont={\small\bfseries\color{blue!70!black}},
   ]%
   \small\ttfamily
  }
  {%
   \end{mdframed}%
  }

% --- Código JavaScript (para fragmentos de app.js) ---
\lstdefinestyle{jsstyle}{
  language      = Java,
  morekeywords  = {const,let,var,module,exports},
  basicstyle    = \small\ttfamily,
  keywordstyle  = \color{blue!70!black},
  stringstyle   = \color{red!60!black},
  commentstyle  = \color{gray!50},
  numbers       = none,
  showstringspaces=false,
  breaklines    = true,
  frame         = none
}

% ---------- Portada ----------
\title{\textbf{Flujo de Programación para Principiantes}\\[6pt]
       \large 10 ejercicios resueltos con diagramas Mermaid}
\author{Basado en \texttt{4GeeksAcademy/program-flow-for-beginners}}
\date{}

\begin{document}
\maketitle

\tableofcontents
\newpage

% ==================================================================
\section{Introducción}
% ==================================================================

Los diagramas de flujo son una herramienta visual fundamental para entender
la lógica de un programa antes de escribirlo. En esta guía resolvemos los
10 ejercicios del curso \emph{Program Flow for Beginners} usando la
notación \textbf{Mermaid}, un lenguaje de marcado que permite describir
diagramas de flujo con texto plano.

Cada ejercicio presenta:
\begin{itemize}[nosep]
  \item El escenario del problema.
  \item Las reglas que debe cumplir el diagrama (rúbrica).
  \item El diagrama Mermaid resuelto.
  \item Una explicación paso a paso.
\end{itemize}

Los archivos \texttt{app.js} contienen el diagrama dentro de una cadena
de texto (\texttt{const answer = `...`}) que se exporta para ser validada
automáticamente.

\newpage

% ==================================================================
\section{Ejercicio 01 — Sequence Basics (Flujo Lineal)}
% ==================================================================

\textbf{Escenario:} Preparas café antes de clases. Define un flujo
lineal: inicio, hervir agua, preparar café, servir, fin. Sin decisiones.

\textbf{Reglas (rúbrica):}
\begin{itemize}
  \item Nodos requeridos: \texttt{start}, \texttt{end}.
  \item Aristas requeridas: \texttt{A→B}, \texttt{B→C}, \texttt{C→D}.
  \item Mínimo 3 aristas.
  \item Prohibido: \texttt{while}, \texttt{for}.
\end{itemize}

\textbf{Código en \texttt{app.js}:}

\begin{lstlisting}[style=jsstyle]
const answer = `
flowchart TD
    A[start] --> B[step 1]
    B --> C[step 2]
    C --> D[end]
`;
module.exports = answer.trim();
\end{lstlisting}

\begin{mermaid}
flowchart TD
    A[start] --> B[step 1]
    B --> C[step 2]
    C --> D[end]
\end{mermaid}

\textbf{Explicación:}
Es el patrón más simple: \textbf{secuencia lineal}. No hay bifurcaciones
ni bucles. Cada paso ocurre uno después del otro, como seguir una
receta de cocina.

\begin{enumerate}
  \item \texttt{A[start]} — punto de entrada del programa.
  \item \texttt{B[step 1]} — primera acción (ej: hervir agua).
  \item \texttt{C[step 2]} — segunda acción (ej: preparar café).
  \item \texttt{D[end]} — fin del flujo.
\end{enumerate}

Los IDs (\texttt{A}, \texttt{B}, \texttt{C}, \texttt{D}) son
obligatorios según la rúbrica. La flecha \texttt{-->} indica la
dirección del flujo.

\newpage

% ==================================================================
\section{Ejercicio 02 — If Else Basics (Decisión Binaria)}
% ==================================================================

\textbf{Escenario:} Estás en una fiesta. Una persona solo puede beber
alcohol si tiene 18 o más. Crea el flujo para validar edad: ruta
permitida, ruta no permitida y fin.

\textbf{Reglas:}
\begin{itemize}
  \item Nodos: \texttt{start}, \texttt{input}, \texttt{valid},
        \texttt{invalid}, \texttt{end}.
  \item Aristas: \texttt{A→B}, \texttt{B→C}, \texttt{C→D},
        \texttt{C→E}, \texttt{D→F}, \texttt{E→F}.
  \item Condiciones: \texttt{yes}, \texttt{no}.
  \item Mínimo 6 aristas.
\end{itemize}

\begin{lstlisting}[style=jsstyle]
const answer = `
flowchart TD
    A[start] --> B[input]
    B --> C{decision}
    C --|yes|--> D[valid]
    C --|no|--> E[invalid]
    D --> F[end]
    E --> F[end]
`;
module.exports = answer.trim();
\end{lstlisting}

\begin{mermaid}
flowchart TD
    A[start] --> B[input]
    B --> C{decision}
    C --|yes|--> D[valid]
    C --|no|--> E[invalid]
    D --> F[end]
    E --> F[end]
\end{mermaid}

\textbf{Explicación:}
Introducimos el primer \textbf{condicional}. El diamante \texttt{C\{decision\}}
representa una pregunta que tiene dos respuestas posibles:
\begin{itemize}
  \item \texttt{yes} (edad ≥ 18) → \texttt{D[valid]} (puede beber).
  \item \texttt{no}  (edad < 18) → \texttt{E[invalid]} (no puede).
\end{itemize}
Ambos caminos convergen al final (\texttt{F[end]}). Esto se llama
\textbf{convergencia} y evita que el diagrama quede con extremos
sueltos.

\newpage

% ==================================================================
\section{Ejercicio 03 — Nested Decisions (Decisiones Anidadas)}
% ==================================================================

\textbf{Escenario:} Alquiler de auto. Solo puedes alquilar si tienes
21+ años \textbf{y} licencia de conducir válida. Modela decisiones
anidadas y resultados finales.

\textbf{Reglas:}
\begin{itemize}
  \item Nodos: \texttt{start}, \texttt{input}, \texttt{valid},
        \texttt{home}, \texttt{end}.
  \item Aristas: \texttt{A→B}, \texttt{B→C}, \texttt{C→D},
        \texttt{D→E}, \texttt{C→F}, \texttt{D→F}, \texttt{E→F}.
  \item Condiciones: \texttt{yes}, \texttt{no}.
  \item Mínimo 7 aristas.
\end{itemize}

\begin{lstlisting}[style=jsstyle]
const answer = `
flowchart TD
    A[start] --> B[input]
    B --> C{age >= 21?}
    C --|yes|--> D{license valid?}
    C --|no|--> F[end]
    D --|yes|--> E[valid]
    D --|no|--> F[end]
    E --> F[end]
`;
module.exports = answer.trim();
\end{lstlisting}

\begin{mermaid}
flowchart TD
    A[start] --> B[input]
    B --> C{age >= 21?}
    C --|yes|--> D{license valid?}
    C --|no|--> F[end]
    D --|yes|--> E[valid]
    D --|no|--> F[end]
    E --> F[end]
\end{mermaid}

\textbf{Explicación:}
Aquí una decisión está \textbf{anidada} dentro de otra. Primero se
pregunta por la edad; solo si es suficiente se pregunta por la licencia.
\begin{enumerate}
  \item \texttt{C\{age >= 21?\}} — Primera condición.
  \item Si es \texttt{no} → termina directamente.
  \item Si es \texttt{yes} → pasa a la segunda condición
        \texttt{D\{license valid?\}}.
  \item Dentro de \texttt{D}: \texttt{yes} → \texttt{valid},
        \texttt{no} → \texttt{end}.
\end{enumerate}
Este patrón modela condiciones compuestas con AND lógico.

\newpage

% ==================================================================
\section{Ejercicio 04 — Multi Branch (Ramas Múltiples)}
% ==================================================================

\textbf{Escenario:} Control de semáforo. Según el color (verde,
amarillo, rojo), elige la acción correspondiente y luego continúa.

\textbf{Reglas:}
\begin{itemize}
  \item Nodos: \texttt{start}, \texttt{input}, \texttt{home},
        \texttt{output}, \texttt{end}.
  \item Aristas: \texttt{A→B}, \texttt{B→C}, \texttt{B→D},
        \texttt{C→E}.
  \item Condiciones: \texttt{yes}, \texttt{no}.
  \item Mínimo 4 aristas.
\end{itemize}

\begin{lstlisting}[style=jsstyle]
const answer = `
flowchart TD
    A[start] --> B[input]
    B --|yes|--> C[output]
    B --|no|--> D[home]
    C --> E[end]
`;
module.exports = answer.trim();
\end{lstlisting}

\begin{mermaid}
flowchart TD
    A[start] --> B[input]
    B --|yes|--> C[output]
    B --|no|--> D[home]
    C --> E[end]
\end{mermaid}

\textbf{Explicación:}
A diferencia del \texttt{if-else} simple, aquí desde \texttt{B[input]}
salen \textbf{dos aristas etiquetadas} simultáneamente. El nodo
\texttt{B} funciona como un \texttt{switch} o decisión de múltiples
ramas:
\begin{itemize}
  \item Rama \texttt{yes} → procesa (\texttt{output}) y luego termina.
  \item Rama \texttt{no} → va a casa (\texttt{home}) — el flujo termina
        ahí sin necesidad de llegar a \texttt{end} (aunque en la rúbrica
        \texttt{home} es el final para esa rama).
\end{itemize}
La arista \texttt{D[home]} queda como nodo terminal sin salida, lo cual
es válido en este patrón.

\newpage

% ==================================================================
\section{Ejercicio 05 — Loops While (Bucle While)}
% ==================================================================

\textbf{Escenario:} Inicio de sesión con contraseña. Sigue pidiendo la
contraseña mientras sea incorrecta; cuando sea correcta, concede acceso
y termina.

\textbf{Reglas:}
\begin{itemize}
  \item Nodos: \texttt{start}, \texttt{input}, \texttt{loop},
        \texttt{end}.
  \item Aristas: \texttt{A→B}, \texttt{B→C}, \texttt{C→B},
        \texttt{C→D}.
  \item Condiciones: \texttt{yes}, \texttt{no}.
  \item Prohibido: \texttt{for}.
  \item Mínimo 4 aristas.
\end{itemize}

\begin{lstlisting}[style=jsstyle]
const answer = `
flowchart TD
    A[start] --> B[input]
    B --> C{loop}
    C --|no|--> B
    C --|yes|--> D[end]
`;
module.exports = answer.trim();
\end{lstlisting}

\begin{mermaid}
flowchart TD
    A[start] --> B[input]
    B --> C{loop}
    C --|no|--> B
    C --|yes|--> D[end]
\end{mermaid}

\textbf{Explicación:}
Este es el primer \textbf{bucle} (ciclo). La flecha \texttt{C --|no|--> B}
crea un retroceso: si la contraseña es incorrecta (\texttt{no}), el
flujo vuelve a pedir input. Esto es la esencia de un bucle
\texttt{while}:

\begin{quote}
Mientras (contraseña incorrecta) \{ pedir de nuevo \}
\end{quote}

\begin{itemize}
  \item \texttt{C{loop}} evalúa la condición.
  \item \texttt{no} → regresa a \texttt{B[input]} (reintentar).
  \item \texttt{yes} → sale del bucle y termina.
\end{itemize}

\newpage

% ==================================================================
\section{Ejercicio 06 — Loops For (Bucle For)}
% ==================================================================

\textbf{Escenario:} Calificar 5 tareas. Itera desde la tarea 1 hasta la
5, procesa cada una, luego muestra resumen y finaliza.

\textbf{Reglas:}
\begin{itemize}
  \item Nodos: \texttt{start}, \texttt{input}, \texttt{loop},
        \texttt{output}, \texttt{end}.
  \item Aristas: \texttt{A→B}, \texttt{B→C}, \texttt{C→D},
        \texttt{D→C}, \texttt{C→E}.
  \item Condiciones: \texttt{yes}, \texttt{no}.
  \item Mínimo 5 aristas.
\end{itemize}

\begin{lstlisting}[style=jsstyle]
const answer = `
flowchart TD
    A[start] --> B[input]
    B --> C{loop}
    C --|yes|--> D[output]
    C --|no|--> E[end]
    D --> C
`;
module.exports = answer.trim();
\end{lstlisting}

\begin{mermaid}
flowchart TD
    A[start] --> B[input]
    B --> C{loop}
    C --|yes|--> D[output]
    C --|no|--> E[end]
    D --> C
\end{mermaid}

\textbf{Explicación:}
Similar al while, pero la condición del bucle se evalúa de forma
distinta:
\begin{itemize}
  \item \texttt{yes} → procesa la tarea actual (\texttt{output}) y
        luego \texttt{D --> C} devuelve el flujo al bucle para la
        siguiente iteración.
  \item \texttt{no} → todas las tareas están procesadas, sale del
        bucle y termina.
\end{itemize}

Este patrón modela un \texttt{for (i = 1; i <= 5; i++)}. La arista
\texttt{D --> C} es la que materializa la repetición.

\newpage

% ==================================================================
\section{Ejercicio 07 — Input Validation (Validación de Entrada)}
% ==================================================================

\textbf{Escenario:} Email de registro de usuario. Solicita email y
valida formato; si es inválido, pedir de nuevo; si es válido, continuar
y terminar.

\textbf{Reglas:}
\begin{itemize}
  \item Nodos: \texttt{start}, \texttt{input}, \texttt{valid},
        \texttt{invalid}, \texttt{end}.
  \item Aristas: \texttt{A→B}, \texttt{B→C}, \texttt{C→D},
        \texttt{D→B}, \texttt{C→E}.
  \item Condiciones: \texttt{valid}, \texttt{invalid}.
  \item Mínimo 5 aristas.
\end{itemize}

\begin{lstlisting}[style=jsstyle]
const answer = `
flowchart TD
    A[start] --> B[input]
    B --> C{valid}
    C --|invalid|--> D[invalid]
    D --> B
    C --|valid|--> E[end]
`;
module.exports = answer.trim();
\end{lstlisting}

\begin{mermaid}
flowchart TD
    A[start] --> B[input]
    B --> C{valid}
    C --|invalid|--> D[invalid]
    D --> B
    C --|valid|--> E[end]
\end{mermaid}

\textbf{Explicación:}
Un patrón muy común en programas reales: \textbf{validación con
reintento}. Las etiquetas de condición ya no son \texttt{yes/no} sino
\texttt{valid} e \texttt{invalid}, más semánticas:
\begin{itemize}
  \item \texttt{invalid} → email mal formado. El flujo va a
        \texttt{D[invalid]} y luego \texttt{D --> B} regresa a pedir
        el email de nuevo.
  \item \texttt{valid} → email correcto. Sale del bucle y termina.
\end{itemize}
La diferencia con el While (ej. 05) es que aquí la condición del bucle
está etiquetada con palabras específicas del dominio (\texttt{valid},
\texttt{invalid}) y no con \texttt{yes/no}.

\newpage

% ==================================================================
\section{Ejercicio 08 — Accumulator Pattern (Patrón Acumulador)}
% ==================================================================

\textbf{Escenario:} Total de carrito de compras. Lee precios
repetidamente, acumula total y cuando el usuario termina, muestra total
y finaliza.

\textbf{Reglas:}
\begin{itemize}
  \item Nodos: \texttt{start}, \texttt{input}, \texttt{loop},
        \texttt{output}, \texttt{end}.
  \item Aristas: \texttt{A→B}, \texttt{B→C}, \texttt{C→B},
        \texttt{C→D}, \texttt{D→E}.
  \item Condiciones: \texttt{yes}, \texttt{no}.
  \item Mínimo 5 aristas.
\end{itemize}

\begin{lstlisting}[style=jsstyle]
const answer = `
flowchart TD
    A[start] --> B[input]
    B --> C{loop}
    C --|no|--> B
    C --|yes|--> D[output]
    D --> E[end]
`;
module.exports = answer.trim();
\end{lstlisting}

\begin{mermaid}
flowchart TD
    A[start] --> B[input]
    B --> C{loop}
    C --|no|--> B
    C --|yes|--> D[output]
    D --> E[end]
\end{mermaid}

\textbf{Explicación:}
El patrón \textbf{acumulador} es la base de operaciones como sumar
precios, contar elementos o concatenar strings. La estructura es la de
un bucle while con un paso adicional de salida:
\begin{enumerate}
  \item \texttt{B[input]} — leer un precio (o la señal de "terminar").
  \item \texttt{C{loop}} — ¿el usuario terminó?
  \item \texttt{no} → volver a \texttt{B} (seguir acumulando).
  \item \texttt{yes} → mostrar el total acumulado (\texttt{output}) y
        terminar.
\end{enumerate}

En código se vería así:
\begin{lstlisting}[style=jsstyle]
let total = 0;
let precio;
while (precio !== "fin") {
  precio = leerPrecio();
  total += precio;
}
mostrarTotal(total);
\end{lstlisting}

\newpage

% ==================================================================
\section{Ejercicio 09 — Stateful Flow (Flujo con Estado)}
% ==================================================================

\textbf{Escenario:} Máquina de estados de torniquete. Empieza bloqueado;
procesa eventos (moneda/empujar), actualiza estado y permite pasar solo
en transiciones válidas.

\textbf{Reglas:}
\begin{itemize}
  \item Nodos: \texttt{start}, \texttt{input}, \texttt{valid},
        \texttt{output}, \texttt{end}.
  \item Aristas: \texttt{A→B}, \texttt{B→C}, \texttt{C→B},
        \texttt{C→D}, \texttt{D→E}.
  \item Condiciones: \texttt{yes}, \texttt{no}.
  \item Mínimo 5 aristas.
\end{itemize}

\begin{lstlisting}[style=jsstyle]
const answer = `
flowchart TD
    A[start] --> B[input]
    B --> C{valid}
    C --|no|--> B
    C --|yes|--> D[output]
    D --> E[end]
`;
module.exports = answer.trim();
\end{lstlisting}

\begin{mermaid}
flowchart TD
    A[start] --> B[input]
    B --> C{valid}
    C --|no|--> B
    C --|yes|--> D[output]
    D --> E[end]
\end{mermaid}

\textbf{Explicación:}
Este diagrama modela una \textbf{máquina de estados} simple. El nodo
\texttt{C{valid}} representa la validación de la transición actual:
\begin{itemize}
  \item \texttt{B[input]} recibe el evento (moneda o empujón).
  \item \texttt{C{valid}} verifica si la transición es válida según
        el estado actual.
  \item \texttt{no} → el evento es inválido (ej: empujar cuando está
        bloqueado), se ignora y se vuelve a esperar input.
  \item \texttt{yes} → transición válida (ej: moneda insertada),
        produce \texttt{output} (abra paso) y termina.
\end{itemize}

Aunque la estructura es similar al acumulador, la semántica es
diferente: aquí \texttt{C} representa una \textbf{guardia} que filtra
eventos basada en el estado interno.

\newpage

% ==================================================================
\section{Ejercicio 10 — Integrated Flow (Flujo Integrado)}
% ==================================================================

\textbf{Escenario:} Máquina expendedora. Valida selección, maneja bucle
de entrada inválida, verifica pago, entrega producto y completa el
flujo.

\textbf{Reglas:}
\begin{itemize}
  \item Nodos: \texttt{start}, \texttt{input}, \texttt{valid},
        \texttt{invalid}, \texttt{home}, \texttt{output}, \texttt{end}.
  \item Aristas: \texttt{A→B}, \texttt{B→C}, \texttt{B→D},
        \texttt{C→E}, \texttt{D→B}, \texttt{E→F}, \texttt{F→G}.
  \item Condiciones: \texttt{yes}, \texttt{no}, \texttt{valid},
        \texttt{invalid}.
  \item Mínimo 7 aristas.
\end{itemize}

\begin{lstlisting}[style=jsstyle]
const answer = `
flowchart TD
    A[start] --> B[input]
    B --|valid|--> C[valid]
    B --|invalid|--> D[invalid]
    D --|no|--> B
    C --|yes|--> E[home]
    E --> F[output]
    F --> G[end]
`;
module.exports = answer.trim();
\end{lstlisting}

\begin{mermaid}
flowchart TD
    A[start] --> B[input]
    B --|valid|--> C[valid]
    B --|invalid|--> D[invalid]
    D --|no|--> B
    C --|yes|--> E[home]
    E --> F[output]
    F --> G[end]
\end{mermaid}

\textbf{Explicación:}
Es el ejercicio más completo porque \textbf{combina} múltiples patrones
en un solo flujo:

\begin{enumerate}
  \item Desde \texttt{B[input]} salen \textbf{dos ramas etiquetadas}
        con \texttt{valid} e \texttt{invalid} (multi-branch, como en
        ej. 04).
  \item \textbf{Rama inválida}: \texttt{D[invalid]} y luego
        \texttt{no} → vuelve a \texttt{B} (bucle de validación,
        como en ej. 05 y 07).
  \item \textbf{Rama válida}: \texttt{C[valid]} → \texttt{yes} →
        \texttt{E[home]} (transición de estado, como en ej. 09).
  \item Luego el flujo continúa linealmente: \texttt{output} →
        \texttt{end} (secuencia, como en ej. 01).
\end{enumerate}

Este diagrama integra: secuencia lineal, decisión binaria, múltiples
ramas, bucle de validación y máquina de estados en un solo flujo
coherente.

\newpage

% ==================================================================
\section{Resumen de Patrones}
% ==================================================================

\begin{center}
\begin{tabular}{|c|c|c|}
\hline
\textbf{Ejercicio} & \textbf{Patrón principal} & \textbf{Elemento nuevo} \\
\hline
01 & Secuencia lineal & Flechas simples \\
02 & If-else & Diamante + yes/no \\
03 & Decisiones anidadas & Diamante dentro de rama \\
04 & Multi-rama & Dos aristas desde un nodo \\
05 & Bucle while & Arista que retrocede \\
06 & Bucle for & Procesar y luego retroceder \\
07 & Validación & Condiciones semánticas \\
08 & Acumulador & Bucle + salida después \\
09 & Máquina de estados & Guardia de eventos \\
10 & Integrado & Todos los patrones juntos \\
\hline
\end{tabular}
\end{center}

\vspace{10pt}

\textbf{Consejos finales:}

\begin{itemize}
  \item Los IDs de nodos (\texttt{A}, \texttt{B}, \texttt{C}, \dots)
        son obligatorios según la rúbrica de cada ejercicio.
  \item Las etiquetas de las aristas se escriben con la sintaxis
        \texttt{--|\emph{etiqueta}|-->}.
  \item Los nodos pueden ser rectangulares (\texttt{[texto]}),
        de decisión (\texttt{\{texto\}}) o redondeados (\texttt{(texto)}).
  \item Para crear un bucle, la arista debe apuntar a un nodo
        \textbf{anterior} en el flujo.
\end{itemize}

\end{document}